\clearpage
\section{Binary semantic weakening with independent operands}
\label{sec:binary-semantic-weakening}
This milestone completes the application and binary-operation part of the
weakening calculation, with arbitrary independent operands and all their
old free variables and atoms retained. It also proves compatibility with
the original maximum-depth padding and exact contextual observations.
Together with Section~\ref{sec:resource-semantic-weakening}, this supplies
the graph calculations for all non-atomic term constructors. The atomic
translation cases and full syntactic induction are completed in
Section~\ref{sec:term-semantic-weakening}.
No term, Figure~16 translation clause, or machine rule has changed.

The paper's weakening proof reduces two unused inputs of a binary distributor
by Idem. Here those inputs come from different operand graphs, so the
previous copying calculation for two copies of one payload does not suffice.
The new proof identifies the two actual weakenings, cuts their actual common
row out of the whole graph, and proves that the remaining graph is exactly
the two old operands with all old distributor rows. The source is
\href{https://arxiv.org/html/1907.01257v5\#S12.SS3}{Section~12.3 of the paper},
the weakening proof preceding Proposition~12.2.

\paragraph{Notation and scope.}
Write $E_{k,h}=*^k\diamond^h$, and let $W_sG$ extend $G$ by one unused
resource of sort $s$. The new resource lies at position $k$: a new variable
is appended to the variable block, or a new atom is prepended to the atom
block. Let $\iota_s$ be the resulting injection of old output positions.
A calculation $G\leadsto H$ is the existing finite proof tree of genuine
Idem and BPullW fillings, ordered recursive isomorphisms, reversal and
composition. It is not an extra execution rule or an assumed equivalence.

\Needspace{12\baselineskip}
\subsection{A complete partition and its actual residual}
\termdef{unused-binary-partition}
Suppose $G:I\to E_{k,h}$ is the old graph and $H:I\to E_{k',h'}$ is a
candidate extended graph. A partition consists of complete-label occurrences
of $G$ and the canonical binary row fed by two weakenings in $H$. Their edge
images must be disjoint and jointly exhaustive. The old occurrence fixes
inputs and maps output $i$ to $\iota_s(i)$; the row's sole output is the new
position $k$. These are geometric data, not semantic assumptions. They are
all constructed below for the original paired translation.
\lean{Copying/WeakeningRowPartition}{13}{27}{Weakening.RowPartition.Partition}

The corrected graph is the computed finite complement with the selected
five-edge pattern replaced by one primitive weakening. The old occurrence
survives this replacement because it is edge-disjoint from the selected row.
The new weakening has its literal occurrence in the inserted component of
the replacement. Boundary transport accounts for the equal hole interfaces;
it does not change a port or an incidence.
\lean{Copying/WeakeningRowPartition}{29}{31}{Weakening.RowPartition.corrected}
\lean{Copying/WeakeningRowPartition}{34}{37}{Weakening.RowPartition.retained}
\lean{Copying/WeakeningRowPartition}{40}{47}{Weakening.RowPartition.inserted}

The two final components are the retained old graph and that inserted
weakening. Their edges have different replacement tags, hence are disjoint.
Their ordered parallel assembly, using the original insertion boundary,
gives an occurrence of $W_sG$ in the corrected graph. Its old-component
port equation follows from the parallel occurrence's embedding equation.
\lean{Copying/WeakeningRowPartition}{48}{54}{Weakening.RowPartition.finalComponents}
\lean{Copying/WeakeningRowPartition}{55}{63}{Weakening.RowPartition.finalComponents_disjoint}
\lean{Copying/WeakeningRowPartition}{66}{73}{Weakening.RowPartition.normalizationOccurrence}
\lean{Copying/WeakeningRowPartition}{74}{80}{Weakening.RowPartition.normalization_old_port}

For an old input or output, the injectivity of the finite-complement embedding
identifies its cut coordinate with the expected external coordinate.
The replacement therefore leaves every old input fixed and every old output
at $\iota_s(i)$. The same argument at the selected row's output places the
inserted weakening at $k$. The reconstruction's input equation follows by
factoring through the old occurrence. Its component-output equation treats
the old and new components separately; the inverse law of the insertion
boundary then proves that every final output position is fixed.
\lean{Copying/WeakeningRowPartition}{81}{93}{Weakening.RowPartition.retained_input}
\lean{Copying/WeakeningRowPartition}{94}{106}{Weakening.RowPartition.retained_output}
\lean{Copying/WeakeningRowPartition}{107}{122}{Weakening.RowPartition.inserted_output}
\lean{Copying/WeakeningRowPartition}{123}{126}{Weakening.RowPartition.normalization_input}
\lean{Copying/WeakeningRowPartition}{127}{149}{Weakening.RowPartition.normalization_component_output}
\lean{Copying/WeakeningRowPartition}{151}{161}{Weakening.RowPartition.normalization_output}

\begin{lemma}[Exhaustive residual reconstruction]
The assembled occurrence of $W_sG$ in the corrected graph is an ordered
recursive isomorphism.
\end{lemma}
\begin{proof}
An edge of the replacement is either inserted or retained. An inserted edge
is the new weakening and is in the assembled image. A retained edge comes
from the complement of the selected row. By the assumed exhaustive partition,
it belongs to the row or to the old occurrence. The first alternative
contradicts complement membership; the second supplies its preimage.
Thus the assembled occurrence is edge-surjective. Together with its complete
incidence laws and fixed external interfaces, the existing occurrence-cover
lemma upgrades it to a recursive isomorphism.
\end{proof}
\lean{Copying/WeakeningRowPartition}{164}{184}{Weakening.RowPartition.normalization_edge_surjective}
\lean{Copying/WeakeningRowPartition}{187}{191}{Weakening.RowPartition.normalizationIso}

\begin{theorem}[Two-Idem partition normalization]
For any such partition, if $H$ is recursively focus-free, then
$H\leadsto W_sG$ by two genuine Idem fillings.
\end{theorem}
\begin{proof}
Apply the previously proved first Idem cut witness at the selected occurrence,
then the second witness in the same computed complement. Compose the second
witness's target representation with the inverse reconstruction isomorphism.
The result is precisely $W_sG$, with every retained old component unchanged.
\end{proof}
\lean{Copying/WeakeningRowPartition}{194}{197}{Weakening.RowPartition.normalize}

\Needspace{12\baselineskip}
\subsection{Keeping two unrelated operands and their old distributor}
Let $G,H:*\to E_{k,h}$ be arbitrary independent graphs. Their pair $P(G,H)$
is the original Figure~16 ordered parallel assembly, which moves the left
atom block past the right variable block. The auxiliary presentation exposes
the output equality as explicit boundary transport. An equality theorem,
not merely an isomorphism, identifies it with the original translation pair.
Define the binary core by
\[
        C(G,H)=P(G,H);D_{k,h,2},
\]
where $D_{k,h,2}$ is the original term-and-atom environment distributor.
\lean{Copying/WeakeningPairOccurrences}{13}{18}{Weakening.Pair.paired}
\lean{Copying/WeakeningPairOccurrences}{19}{21}{Weakening.Pair.original_pair}
\lean{Copying/WeakeningPairOccurrences}{23}{26}{Weakening.Pair.core}

Either operand occurs in the pair with its own Boolean component tag.
The input equation places the left and right inputs at positions $0$ and $1$.
The output equation uses the existing pair-boundary index: old variable $i$
at use $b$ has position $i+kb$, and atom positions follow the complete term
block. Both equations follow directly from the prescribed shuffle.
\lean{Copying/WeakeningPairOccurrences}{28}{33}{Weakening.Pair.component}
\lean{Copying/WeakeningPairOccurrences}{34}{37}{Weakening.Pair.component_output}
\lean{Copying/WeakeningPairOccurrences}{38}{46}{Weakening.Pair.component_input}

Compose the old payload occurrence in $W_sG$ or $W_sH$ with the corresponding
component occurrence in their enlarged pair. These two occurrences are
disjoint by their Boolean tags, so their ordered parallel assembly retains
the complete old pair. Following any component port through that assembly
agrees with its direct old-component map. In particular, an old output at
pair position $(b,i)$ goes to $(b,\iota_s(i))$, and the two inputs stay fixed.
No equality between $G$ and $H$ is used.
\lean{Copying/WeakeningPairOccurrences}{48}{53}{Weakening.Pair.oldComponent}
\lean{Copying/WeakeningPairOccurrences}{54}{62}{Weakening.Pair.components_disjoint}
\lean{Copying/WeakeningPairOccurrences}{64}{68}{Weakening.Pair.oldPair}
\lean{Copying/WeakeningPairOccurrences}{69}{76}{Weakening.Pair.oldPair_component}
\lean{Copying/WeakeningPairOccurrences}{77}{89}{Weakening.Pair.oldPair_output}
\lean{Copying/WeakeningPairOccurrences}{90}{106}{Weakening.Pair.oldPair_input}

The retained left occurrence is this old pair embedded into the enlarged
core. The retained right occurrence is the already proved old two-use
distributor embedding followed by its serial embedding. Their shared-port
maps agree: every pair-boundary index has the form $(b,i)$, and both sides
send it to $(b,\iota_s(i))$. Serial tags make their edges disjoint.
Serial assembly therefore gives the complete old core occurrence. Its inputs
are fixed and its outputs follow $\iota_s$, by the component input and
distributor output equations.
\lean{Copying/WeakeningPairOccurrences}{108}{110}{Weakening.Pair.oldLeft}
\lean{Copying/WeakeningPairOccurrences}{111}{114}{Weakening.Pair.oldRight}
\lean{Copying/WeakeningPairOccurrences}{115}{122}{Weakening.Pair.old_boundary}
\lean{Copying/WeakeningPairOccurrences}{123}{125}{Weakening.Pair.old_disjoint}
\lean{Copying/WeakeningPairOccurrences}{127}{130}{Weakening.Pair.old}
\lean{Copying/WeakeningPairOccurrences}{131}{138}{Weakening.Pair.old_input}
\lean{Copying/WeakeningPairOccurrences}{139}{145}{Weakening.Pair.old_output}

\Needspace{12\baselineskip}
\subsection{The actual selected five-edge pattern}
\termdef{independent-operand-row}
The selected row is the new resource's literal canonical two-use row inside
the enlarged distributor. One new weakening edge comes from $W_sG$ and
one from $W_sH$. Their operand tags are different, proving injectivity;
the row is in the distributor side of the serial composite, proving
disjointness from both. Each weakening has its original weakening label
and an empty source list.
\lean{Copying/WeakeningBinaryRow}{12}{14}{Weakening.Pair.row}
\lean{Copying/WeakeningBinaryRow}{16}{22}{Weakening.Pair.weakEdge}
\lean{Copying/WeakeningBinaryRow}{23}{25}{Weakening.Pair.weak_injective}
\lean{Copying/WeakeningBinaryRow}{26}{29}{Weakening.Pair.weak_disjoint}
\lean{Copying/WeakeningBinaryRow}{30}{33}{Weakening.Pair.weak_labels}
\lean{Copying/WeakeningBinaryRow}{34}{44}{Weakening.Pair.weak_sources}

For the target equation, use the occurrence's target preservation: the
weakening in either extended operand targets its new external output.
The pair-boundary equation sends this output to $(b,k)$, and serial gluing
makes it the internal shared wire at that position. The selected distributor
row's input equation identifies this very wire with its input $b$. Thus
both added weakenings actually feed that row, with the original order.
The auxiliary target lemma proves this by mapping the singleton target
list through the occurrence; it does not infer connectivity from edge labels.
\lean{Copying/WeakeningBinaryRow}{47}{59}{Weakening.Pair.weak_target_link}
\lean{Copying/WeakeningBinaryRow}{60}{67}{Weakening.Pair.weak_targets}

The established attachment construction now assembles the complete canonical
pattern: two incoming weakenings, together with the row's own canonical
weakening and two contractions. Its sole output is the new resource's
external position $k$, obtained by specializing the selected row output
through the serial right embedding.
\lean{Copying/WeakeningBinaryRow}{69}{72}{Weakening.Pair.unusedRow}
\lean{Copying/WeakeningBinaryRow}{73}{87}{Weakening.Pair.unusedRow_output}

\begin{lemma}[The actual core has the required exhaustive partition]
In $C(W_sG,W_sH)$, the selected five-edge pattern and the retained old core
$C(G,H)$ are edge-disjoint and jointly exhaustive.
\end{lemma}
\begin{proof}
Split a host edge by its outer serial tag. In the operand side, split again
by operand and by the extension tag: it is either an old edge of that
particular operand or its unique new weakening. In the distributor side,
split by sharing sort and resource index. Under variable extension, an
index is either below $k$ or equal to the new last index $k$; all atom rows
are retained. Under atom extension, the new index is $0$ and every positive
index comes from its predecessor; all variable rows are retained. The
within-row edge index stays unchanged in every case. The same distinctions
prove disjointness: new and old operand tags differ, and selected and old
resource indices cannot coincide. This includes empty old resource blocks.
\end{proof}
\lean{Copying/WeakeningBinaryPartition}{11}{32}{Weakening.Pair.row_old_disjoint}
\lean{Copying/WeakeningBinaryPartition}{35}{87}{Weakening.Pair.edge_cover}

Package these proved facts and the previously proved boundary equations as
the required partition. Focus freedom of the core follows from that of
each operand and the language-label invariant of the original distributor.
Focus freedom of each extension adds only an actual weakening edge. The
general two-Idem theorem therefore gives the unconditional geometric law
\[
          C(W_sG,W_sH)\leadsto W_s C(G,H)
\]
for every recursively focus-free $G,H$. All partition and residual premises
have been discharged for these actual endpoints.
\lean{Copying/WeakeningBinaryPartition}{90}{98}{Weakening.Pair.partition}
\lean{Copying/WeakeningBinaryPartition}{100}{103}{Weakening.Pair.core_free}
\lean{Copying/WeakeningBinaryPartition}{106}{111}{Weakening.Pair.normalize}

\Needspace{12\baselineskip}
\subsection{Original application and every binary operation}
\termdef{binary-translation-weakening}
Recursive isomorphism commutes with extension: place the given isomorphism
beside the identity of the new weakening, preserve the insertion boundary,
and transport the equal external interface. This gives the extension
congruence used in the following reassociation.
\lean{Copying/WeakeningBinarySteps}{11}{16}{Weakening.extend_congr}

Let $Q:I\to[*,*]$ be any unchanged focus-free head. The original clause has
parenthesization $(Q;P(G,H));D_{k,h,2}$. Reassociate it to $Q;C(G,H)$.
Apply the complete core law under the unchanged prefix $Q$, move the
outside weakening past that prefix using the proved serial-extension
isomorphism, then reassociate inside the extension. Every isomorphism fixes
the ordered external boundary. This proves
\[
 (Q;P(W_sG,W_sH));D_{k',h',2}
 \leadsto W_s\bigl((Q;P(G,H));D_{k,h,2}\bigr).
\]
\lean{Copying/WeakeningBinarySteps}{19}{36}{Weakening.binary_head_extend}

For application take $Q$ to be the original application edge. For a binary
operation take its original operation edge. Primitive labels establish
focus freedom without an additional premise. Thus the endpoint laws cover
application and all four binary operations, including assignment and atom
equality, with arbitrary independent operands. The graphs retain the
original full distributor, rather than replacing it by ideal wires.
\lean{Copying/WeakeningBinarySteps}{38}{45}{Weakening.application_extend}
\lean{Copying/WeakeningBinarySteps}{48}{56}{Weakening.binary_extend}

Calculations lift through either side of a parallel assembly with any
prescribed boundary. The proof is induction on the actual calculation:
isomorphisms use parallel congruence; an Idem or BPullW leaf uses the genuine
parallel filling certificate with a reflexive certificate for the unchanged
component; reversal and composition lift recursively. Applying left and
right lifting in succession, followed by the pair's equality transport,
gives pair congruence. Adding the unchanged operation prefix and distributor
suffix gives application and binary congruence.
\lean{Copying/WeakeningBinarySteps}{59}{73}{Weakening.Calculation.parallelRight}
\lean{Copying/WeakeningBinarySteps}{74}{88}{Weakening.Calculation.parallelLeft}
\lean{Copying/WeakeningBinarySteps}{91}{98}{Weakening.Calculation.pair}
\lean{Copying/WeakeningBinarySteps}{99}{104}{Weakening.Calculation.application}
\lean{Copying/WeakeningBinarySteps}{105}{111}{Weakening.Calculation.binary}

Consequently, actual subterm calculations $G'\leadsto W_sG$ and
$H'\leadsto W_sH$ imply the corresponding application or binary calculation
with target $W_s$ of the original compound clause. First lift both subterm
calculations, then use the endpoint law. Focus freedom of $G$ and $H$ is
recovered from the supplied calculations and the extension invariant.
These are induction rules consuming proved subterm calculations; they do
not assert the still-open atomic cases or assume the desired compound
observational equivalence.
\lean{Copying/WeakeningBinarySteps}{113}{120}{Weakening.application_step}
\lean{Copying/WeakeningBinarySteps}{123}{130}{Weakening.binary_step}

\Needspace{12\baselineskip}
\subsection{Original common-depth padding and observations}
\termdef{weakening-padding-compatibility}
Padding $W_sG$ by $n$ layers equals extending the graph padded by $n$ layers.
Induction on $n$ uses the existing exact one-layer extension equation.
Transporting an equality of depth indices also commutes with extension,
by elimination of that equality. Their composition proves the same equality
for padding to any common allowance $N\geq d$.
\lean{Copying/WeakeningPadding}{11}{18}{Weakening.extend_pad}
\lean{Copying/WeakeningPadding}{20}{25}{Weakening.extend_castDepth}
\lean{Copying/WeakeningPadding}{27}{32}{Weakening.extend_padTo}

Padding a weakening calculation transports each original Idem or BPullW
leaf by its proved padding-closure law. Rewriting the target by the exact
extension equation gives a padded subterm calculation. For application or
a binary operation with operand depths $d,e$, choose $N=\max(d,e)$, pad
the two independent calculations to $N$, then apply the corresponding
same-depth induction step. This is precisely the original Figure~16
maximum-depth convention, including unequal depths and nested boxes.
\lean{Copying/WeakeningPadding}{35}{42}{Weakening.Calculation.weakening_padTo}
\lean{Copying/WeakeningPadding}{45}{58}{Weakening.application_step_padded}
\lean{Copying/WeakeningPadding}{60}{73}{Weakening.binary_step_padded}

The addressed theorems place each complete endpoint calculation at any
recursive site of a focus-free closing host and apply the already proved
exact-count soundness theorem. Hole freedom of the filled source supplies
hole freedom of intermediates and the target. Both refinements then hold
in the original unbounded machine with equal successful-transition counts.
There is no operand termination premise, and the structural result covers
all intended binding-free observation contexts, including observations
that reuse a closure.
\lean{Copying/WeakeningPadding}{76}{86}{Weakening.binary_addressed}
\lean{Copying/WeakeningPadding}{89}{99}{Weakening.application_addressed}

\Needspace{12\baselineskip}
\subsection{Checked composite instances and dependency map}
The first open body is the direct translation of $f(!a_1)$ with three
variable outputs and two atom outputs. The separate literal graph is the
translation of $1$ with two variable and two atom outputs. These definitions
use the unchanged syntax and translation.
\lean{Copying/WeakeningBinaryExamples}{11}{13}{Weakening.Examples.callbackBody}
\lean{Copying/WeakeningBinaryExamples}{15}{20}{Weakening.Examples.allocationValue}

The assignment instance combines an allocation weakening calculation at
depth zero with an abstraction weakening calculation at depth one, then
uses the padded binary induction step. Its target is exactly the original
translation of $(\mathtt{ref}\ 1):=(\lambda f.\,f(!a_1))$, extended by one
unused resource. Its source exposes the new weakening separately inside
the two operand constructions; it is not asserted to be the translation
of a syntactically weakened term. That additional identification needs the
atomic translation base cases and full induction.
\lean{Copying/WeakeningBinaryExamples}{21}{32}{Weakening.Examples.assignment_closure}

A second body is the translation of $\lambda z.\,f\,z$ and already contains
a box. Abstracting it gives a depth-two operand, independent of the depth-one
callback just described. The padded application calculation has as target
the actual translation of
$(\lambda f.\lambda z.\,f\,z)(\lambda f.\,f(!a_1))$, extended by the unused
resource. Both examples are checked for either sharing sort in a mixed
nonempty environment. A separate finite-carrier calculation verifies that
the removed canonical pattern has five edges and its residual has one,
at every depth allowance. These are structural calculations, not claims
that the displayed open programs terminate.
\lean{Copying/WeakeningBinaryExamples}{34}{36}{Weakening.Examples.nestedCallbackBody}
\lean{Copying/WeakeningBinaryExamples}{39}{50}{Weakening.Examples.unequal_depth_application}
\lean{Copying/WeakeningBinaryExamples}{53}{57}{Weakening.Examples.removed_row_size}

\paragraph{Dependency map and remaining work.}
The completed path is: independent operand occurrences and old distributor
embedding $\to$ actual attached five-edge row $\to$ exhaustive whole-core
partition $\to$ two original Idem fillings and residual isomorphism
$\to$ application and full binary-alphabet endpoint laws $\to$ separate
subterm induction steps $\to$ original maximum-depth padding $\to$ exact
addressed observations. The only semantic leaves are previously proved
Idem and BPullW results. Recursive focus freedom is an application condition,
not an unproved template or machine law.

The atomic base cases, binder-index identifications, full syntactic induction
and arbitrary finite environment-extension calculations are now completed in
Section~\ref{sec:term-semantic-weakening}. Figure~22 and the correspondence and
substitution obligations in Appendix~B remain downstream. Proposition~12.2
is not yet proved.
